\documentclass[../../../include/open-logic-section]{subfiles}

\begin{document}

\olfileid[fr]{sth}{z}{infinity-again}	
\olsection{Infini}

Nous disposons déjà d'assez d'axiomes pour assurer qu'il existe une
infinité d'ensembles (s'il en existe au moins un). Supposons en effet
qu'un ensemble existe : alors $\emptyset$ existe (par la
\olref[sth][z][sep]{prop:emptyexists}). Or, pour tout ensemble~$x$,
l'ensemble $x \cup \{x\}$ existe par la
\olref[sth][z][pairs]{prop:pairsconsequences}. En répétant cette
opération quelques fois, nous obtenons les ensembles suivants :
\begin{enumerate}
	\item[0.] $\emptyset$ %& $0$ & étape $0$\\
	\item[1.] $ \{\emptyset\}$ %& $1$ & étape $1$\\
	\item[2.] $\{\emptyset, \{\emptyset\}\}$ %& $2$ & étape $2$\\
	\item[3.] $\{\emptyset, \{\emptyset\}, \{\emptyset, \{\emptyset\}\}\}$ %&$3$ & étape $3$\\
	\item[4.] $\{\emptyset, \{\emptyset\}, \{\emptyset, \{\emptyset\}\}, \{\emptyset, \{\emptyset\}, \{\emptyset, \{\emptyset\}\}\}\}$% & $4$ & étape $4$
\end{enumerate}
On peut vérifier que ces ensembles sont tous distincts.

Nous avons commencé la numérotation à $0$ pour plusieurs raisons.
En voici une : l'ensemble auquel nous avons attribué l'étiquette
«~$n$~» possède exactement $n$ !!{element}s et se forme,
intuitivement, à la $n$-ième étape.

Mais voilà : nous obtenons ainsi une \emph{infinité d'ensembles}, sans
pour autant garantir l'existence d'un \emph{ensemble infini}, c'est-à-dire
d'un ensemble qui possède une infinité d'!!{element}s. La distinction
est essentielle : sans ensemble infini (au sens de Dedekind), nous ne
pouvons pas construire d'algèbre de Dedekind. Or nous voulons une
telle algèbre pour la considérer comme l'ensemble des nombres naturels.
(Voir la \olref[sfr][infinite][dedekindsproof]{sec}.)

Les axiomes posés jusqu'ici ne garantissent \emph{pas} l'existence
d'un ensemble infini. Nous devons donc ajouter un axiome :

\begin{axiom}[Infini]
Il existe un ensemble $I$ tel que $\emptyset \in I$ et que
$x \cup \{x\} \in I$ chaque fois que $x \in I$.
\begin{align*}
	\exists I( & (\exists o \in I)\forall x\ x \notin o \land {}\\
	& (\forall x \in I)(\exists s \in I)\forall z(z \in s \liff (z \in x \lor z = x)))
\end{align*}
\end{axiom}

L'ensemble $I$ fourni par l'axiome de l'infini est infini au sens de
Dedekind. Soit $s(x)=x\cup\{x\}$ et soit $\omega$ la plus petite partie
de $I$ qui contient $\emptyset$ et qui est stable par $s$.
L'ensemble vide est transitif. Si $x$ est transitif, tout élément de
$s(x)$ est inclus dans $x$, puisqu'il appartient à $x$ ou est égal à
$x$. Donc $s(x)$ est transitif et, par récurrence, tout élément de
$\omega$ l'est. Pour $x\in\omega$, on a $\bigcup s(x)=x$ puisque
$\bigcup x\subseteq x$. Ainsi $s$ est injective sur $\omega$. De plus,
$x\in s(x)$, donc son image omet $\emptyset$.
L'application égale à $s$ sur $\omega$ et à l'identité sur
$I\setminus\omega$ est alors une injection de $I$ dans lui-même dont
l'image omet $\emptyset$.

Le \olref[sfr][infinite][dedekind]{thm:DedekindInfiniteAlgebra} donne alors une algèbre de Dedekind.
Nous prenons pour ensemble des nombres naturels $\omega$, avec
$\emptyset$ comme élément initial et la restriction de $s$ comme
successeur.
\begin{defn}\ollabel{defnomega}
Soit $I$ un ensemble quelconque fourni par l'axiome de l'infini. Soit
$s$ la fonction $s(x) = x \cup \{x\}$. Posons $\omega =
\closureofunder{s}{\emptyset}$. Nous appelons les !!{element}s de
$\omega$ les \emph{nombres naturels} et disons que $n$ résulte de
$n$ applications de $s$ à $\emptyset$.
\end{defn}

Nous pouvons maintenant revenir aux exemples précédents et vérifier
que l'ensemble étiqueté «~$n$~» y sera considéré \emph{comme} le
nombre $n$.

Nous discuterons la portée de cette convention à la
\olref[sth][z][nat]{sec}. Pour l'instant, elle nous permet de
démontrer un résultat intuitif :

\begin{prop}\ollabel{naturalnumbersarentinfinite}
Aucun nombre naturel n'est infini au sens de Dedekind.
\end{prop}

\begin{proof}
La démonstration procède par récurrence, selon la
\olref[sfr][infinite][induction]{thm:dedinfiniteinduction}. Il est
clair que $0 = \emptyset$ n'est pas infini au sens de Dedekind. Pour
l'étape de récurrence, établissons la contraposée : si, par impossible,
$s(n)$ est infini au sens de Dedekind, alors $n$ l'est aussi.

Supposons donc que $s(n)$ soit infini au sens de Dedekind : il existe
une !!{injection} $f$ telle que
$\ran{f}\subsetneq \dom{f} = s(n) = n \cup \{n\}$. Il y a deux cas.

\emph{Premier cas : $n \notin \ran{f}$.} Alors $\ran{f} \subseteq n$
et $f(n) \in n$. Posons $g = \funrestrictionto{f}{n}$ ; nous avons
$\ran{g} = \ran{f} \setminus \{f(n)\} \subsetneq n = \dom{g}$.
Donc $n$ est infini au sens de Dedekind.

\emph{Second cas : $n \in \ran{f}$.} Choisissons
$m \in \dom{f} \setminus \ran{f}$ et définissons une fonction $h$
de domaine $s(n) = n \cup \{n\}$ par :
\[
h(x) =
	\begin{cases}
		f(x) & \text{si }f(x) \neq n\\
		m & \text{si }f(x)=n
	\end{cases}
\]
Les fonctions $h$ et $f$ coïncident partout sauf au point où
$h(f^{-1}(n)) = m \neq n = f(f^{-1}(n))$. Comme $m$ n'appartient pas
à l'image de l'!!{injection} $f$, l'application $h$ est injective.
Son image omet $n$, donc $\ran{h} \subsetneq \dom{h} = s(n)$.
L'argument du premier cas montre alors que $n$ est infini au sens
de Dedekind.
\end{proof}

Reste cependant à savoir comment \emph{justifier} l'axiome de l'infini.
La réponse courte est qu'il faut ajouter un autre principe au récit
que nous avons présenté. Le voici :
\begin{enumerate}
	\item[] \stagesinf. Il existe une étape infinie, c'est-à-dire une
	étape qui (a) n'est pas la première, (b) est précédée d'autres
	étapes, mais (c) n'a pas de prédécesseur immédiat.
\end{enumerate}
L'axiome de l'infini découle directement de ce principe. Nous savons
que le nombre naturel $n$ se forme à l'étape $n$. L'ensemble $\omega$
se forme donc à la première étape infinie. Et $\omega$ lui-même
témoigne de l'axiome de l'infini.

Cela ne fait pourtant que reporter la question : comment justifier
\stagesinf{} ? Comme \stagessucc, ce principe ne faisait pas
explicitement partie de notre récit de la hiérarchie cumulative et
itérative. Mais il y a plus : rien, dans l'idée même d'une hiérarchie
itérative où les ensembles se forment étape par étape, ne nous oblige
à penser que ce processus comporte une étape \emph{infinie}. On peut
concevoir sans contradiction apparente que les étapes soient ordonnées
comme les nombres naturels.

Cela soulève toutefois un problème évident. À la
\olref[sfr][infinite][dedekindsproof]{sec}, nous avons examiné la
«~preuve~» de Dedekind selon laquelle il existe un ensemble infini au
sens de Dedekind (composé de pensées). Elle ne vous a peut-être pas
paru très convaincante. Mais si \stagesinf{} ne nous est pas «~imposé~»
par la conception itérative des ensembles (ou par les «~lois de la
pensée~»), nous restons sans justification intrinsèque de l'existence
d'un ensemble infini au sens de Dedekind.

Il y aurait, bien sûr, beaucoup d'autres choses à dire. Mais nous
espérons que vous pouvez maintenant commencer à réfléchir à ce qu'il
\emph{faudrait} pour justifier un axiome (ou un principe). Dans la
suite, nous admettrons simplement \stagesinf{}.

\end{document}
